\chapter{Beta-Thalassemia Genetic Testing}

\section*{What Is This Test?}
Beta-thalassemia genetic testing identifies pathogenic variants in the HBB gene that reduce or abolish beta-globin production. It complements CBC indices, iron studies, and hemoglobin analysis, particularly when the phenotype is unclear.

\section*{Alternative Names}
HBB Gene Testing; Beta-Globin Gene Sequencing; Beta-Thalassemia Molecular Test.

\section*{Why This Test Is Ordered}
Testing is used when microcytosis or anemia suggests thalassemia, when hemoglobin electrophoresis is inconclusive, for prenatal or preconception counseling, and to clarify carrier status or disease severity in an affected family.

\section*{How This Test Is Performed}
DNA from blood is tested for common HBB variants and, when needed, by sequencing and deletion/duplication analysis. Results are interpreted with HbA2, HbF, CBC indices, iron status, and family history.

\section*{Preparation, Risks, and Limitations}
No fasting is required. Risks are limited to blood collection. A negative targeted panel may miss rare variants, and a genetic result should be interpreted by a laboratory professional or genetic counselor rather than used alone to predict severity.

\section*{Understanding Results}
One pathogenic variant usually indicates beta-thalassemia trait, while two clinically significant variants can cause intermedia or transfusion-dependent disease. The exact variants and their combination determine phenotype and reproductive risk.

\section*{References}
Thalassaemia International Federation guidance on beta-thalassemia diagnosis and genetic counseling.

\clearpage
\section*{Understanding Results and Follow-up}
The laboratory report should identify the specimen, method, units, reference interval or decision limit, and whether the result is positive, negative, abnormal, or inconclusive. Reference intervals vary with age, sex, pregnancy, specimen type, assay, and laboratory; a result outside a range is not automatically a diagnosis, and a result within a range does not always exclude disease. Discuss unexpected results with the ordering clinician, who can decide whether repeat or confirmatory testing, treatment, or referral is appropriate. Do not change medicines or delay urgent care based on an isolated result.


\section*{Regulatory Status and Authorization Date}
Regulatory status applies to the specific assay, instrument, manufacturer, and intended use rather than to the test name alone. No single approval date can be assigned to this test category. For a commercial assay, verify the current product in the relevant regulator's database: the U.S. Food and Drug Administration (FDA) clearance, approval, or authorization; India's Central Drugs Standard Control Organisation (CDSCO) licence or permission; the European Union In Vitro Diagnostic Regulation (EU IVDR) CE certificate issued through the applicable conformity-assessment route; or the United Kingdom Medicines and Healthcare products Regulatory Agency (MHRA) registration and UKCA/CE status. Laboratory-developed tests may not have an individual FDA, CDSCO, EU IVDR, or MHRA approval date.
\clearpage
